% Section 5. Drafted from docs/PAPER-OUTLINE.md section 5 and docs/THEORY.md "Layer A". Every named
% result is a gated, axiom-free theorem in coq/*.v (PresentedExecution.v unless noted).

\section{Layer A: execution}\label{sec:layer-a}

Layer~A asks which executions count as the same. Every delivery model in the regime map answers that
question differently: free delivery allows every reordering, causal delivery only those consistent
with happens-before, at-least-once delivery also identifies a duplicate with its first landing, and a
live reconfiguration identifies an event before the switch with its translation after it. This section
states the question once, in a form every model fits, and proves the two principles that answer it.
The exact history results of the map are then instances, recovered with their original statements.

% ---------------------------------------------------------------------------------------------
\subsection{Presented executions}

Executions are words over a set of labels $L$, but not every word is an execution: causal
consistency, typed phases and legal duplicate landings rule some out. The data of a delivery model is
therefore a \emph{presented admissible path system}:
\begin{itemize}[nosep]
  \item a predicate $\mathrm{Adm}$ on words, the admissible executions, closed under prefixes where a
    theorem needs it ($\mathrm{AdmPrefix}$);
  \item generator pairs $\mathrm{Gen}(p, u, v)$: at the point $p$, the fragments $u$ and $v$ are
    interchangeable, for example a swap of independent events $[a; b] \sim [b; a]$, a legal duplicate
    landing $[a; a] \sim [a]$, or a switch square between typed phases;
  \item the generated equivalence $\mathrm{PEq}$: the equivalence closure, through admissible words
    only, of the whiskered pairs $p \cdot u \cdot w \sim p \cdot v \cdot w$.
\end{itemize}
Free delivery is the degenerate case: with every word admissible and generators independent of the
point, $\mathrm{PEq}$ is the monoid congruence the pairs generate (\coq{free_congruence}), and for
adjacent swaps it is equality up to permutation (\coq{swap_peq_perm}). This is the quotient a string
presentation defines, as in trace theory~\cite{cartier1969,mazurkiewicz1977}.

A semantics $\mathrm{sem}$ on words is compared by an equivalence $\approx$; for a state machine,
$\mathrm{sem}(t)$ is the repaired state after running $t$ from the start. \emph{Invariance}
($\mathrm{PInv}$) says that $\mathrm{sem}$ is constant up to $\approx$ on each $\mathrm{PEq}$ class.
This is history descent (\cref{def:history}) for the delivery model the presentation describes.

% ---------------------------------------------------------------------------------------------
\subsection{The two principles}

The \emph{continuation closure} $\mathrm{PPC}$ is the least relation that contains every seed
$(p \cdot u,\, p \cdot v)$ with both sides admissible and is closed under appending the same event to
both sides while both stay admissible. Under prefix closure it is exactly the set of admissible
whiskered pairs (\coq{ppc_iff_whisk}).

\begin{theorem}[The closure principle; \coq{closure_exact}]\label{thm:p2}
If $\mathrm{Adm}$ is closed under prefixes, then invariance holds if and only if $\mathrm{sem}$ agrees,
up to $\approx$, on every pair in the continuation closure.
\end{theorem}

Invariance always implies agreement on the seeds themselves (\coq{seeds_necessary}), but the seeds
alone are not enough. A hidden register makes the gap concrete (\coq{hw_seeds_insufficient}): events
write a register the observation does not show, and a later event reveals it. Every swap seed is then
observationally equal at its point, yet invariance fails; the separating pair is a seed continued by
the reveal (\coq{hw_closure_witness}). The next principle says when the seeds do suffice.

\begin{theorem}[The local principle; \coq{local_exact}]\label{thm:p1}
Let $Q$ be a predicate on words that holds on both sides of every seed and is kept by admissible
extension, and suppose the semantics respects the future on $Q$: whenever $t$ and $t'$ satisfy $Q$,
$t \cdot e$ and $t' \cdot e$ are admissible, and $\mathrm{sem}(t) \approx \mathrm{sem}(t')$, then
$\mathrm{sem}(t \cdot e) \approx \mathrm{sem}(t' \cdot e)$. If $\mathrm{Adm}$ is closed under
prefixes, invariance holds if and only if every generator pair is observationally equal at its
point: $\mathrm{sem}(p \cdot u) \approx \mathrm{sem}(p \cdot v)$ whenever $\mathrm{Gen}(p, u, v)$ and both
sides are admissible.
\end{theorem}

With $Q$ true everywhere this is \coq{local_exact_all}, and the future-respecting hypothesis follows
when the observation is compatible with the steps (\coq{resp_of_compat}). In the hidden register it is
exactly that hypothesis that fails (\coq{hw_not_resp}). Both principles are elementary facts about
congruences generated by local pairs, in the style of semigroup presentations~\cite{howie1995}; the
work is in fitting every delivery model to them.

State descent is not a generator. It runs one word from two states, $s$ and $N(s)$, so it is a
relation on states, not a pair of parallel paths from one point. It enters through the semantics:
under state descent the repaired raw semantics and the governed semantics from $N(s_0)$ are the same
function on words, so they pose the same invariance question for every presentation
(\coq{state_side_transfer}; \cref{lem:transfer} is the special case of equality on histories).
Paths quotient histories; states quotient by repair.

% ---------------------------------------------------------------------------------------------
\subsection{The history results as instances}

Each exact history result of the map is recovered from the local or the closure principle with its original statement, over the
original module's own predicates (\cref{tab:layer-a}).

\begin{table}[t]
\centering\footnotesize
\setlength{\tabcolsep}{4pt}
\begin{tabular}{@{}p{2.4cm}p{4.0cm}p{1.9cm}p{6.0cm}@{}}
\toprule
Delivery model & Recovered as & Principle & Admissible words; generators \\
\midrule
Causal delivery & \coq{PresentedExecution.causal_exact_recovered} & local & causally consistent words; adjacent swaps of concurrent events at a causal point \\
Declared independence & \coq{tconv_exact_recovered} & local & causal histories; swaps of declared pairs (the trace monoid) \\
At-least-once & \coq{causal_alo_idem_recovered}, \coq{alo_exact_recovered} & local & at-least-once deliveries, causal (first name) or in any order (second); concurrent swaps and the landing $[a; a] \sim [a]$ after a first delivery \\
Live switch & \coq{det_live_recovered} & local, $Q$ = past the switch & typed phases; switch squares and post-switch swaps \\
Barrier, modulo a migration & \coq{amodm_closure_recovered} & closure & every word; adjacent swaps (the migration is not compatible with the steps) \\
\bottomrule
\end{tabular}
\caption{Exact history results recovered from the two principles. Each recovered theorem states the
original's conclusion over the original's predicates; the proofs go through \cref{thm:p1,thm:p2}.}
\label{tab:layer-a}
\end{table}

Three of these are not trace monoids. At-least-once delivery identifies a word with a shorter one,
the live switch has typed phases with labels that change at the switch, and the barrier case needs the closure principle
because the migration does not respect the future of the old configuration's steps. The live switch
is the instance where $Q$ matters: the migration does not respect the old configuration's future, but
every seed ends past the switch, where the new configuration's steps respect its observation.

\paragraph{The cost of fitting a model.}
The generic theorems are short; fitting a delivery model to them is not. Each recovery has to show
that the model's notion of ``same delivery'' is exactly the generated equivalence of its presentation,
an adequacy argument about words that does not involve the semantics. For at-least-once delivery that
argument is about 250 lines of word combinatorics. The recoveries therefore compress the history
results into one theorem pair plus a per-model adequacy cost, and the cost is part of the claim.

\paragraph{Not an instance.}
The single-registry theorem's commutation condition CC1 is not a row of this table. Its ambiguity is a
peak of a rewrite system in which repair steps interleave with events, not a pair of event words.
\cref{sec:base-case} combines the local principle on free delivery with a classification of those peaks to derive it.
